%% Adapted from arxiv style template by R Roux 31/01/2019%%
%% Modified by N Hale 09/12/2025%%

%%%%%%%%%%%%%%%%%%%%%%%%%%%%%%%%%%%%%%%%%%%%%%%%%%
%%%%%%% DO NOT CHANGE ANYTHING IN THE PREAMBLE, 
%%%%%%% OTHER THAN ADDING ADDITIONAL PACKAGES/MACROS
%%%%%%%%%%%%%%%%%%%%%%%%%%%%%%%%%%%%%%%%%%%%%%%%%%

\documentclass[12pt,a4paper]{article}

\usepackage[utf8]{inputenc} % allow utf-8 input
\usepackage[T1]{fontenc}    % use 8-bit T1 fonts
\usepackage{xcolor}
\usepackage{hyperref}       % remove comment to add hyperlinks
\hypersetup{
	colorlinks=true,
	linkcolor=blue,
	filecolor=blue,    
	citecolor=blue,  
	urlcolor=blue,
	pdftitle={},
	pdfpagemode=FullScreen,
}
\usepackage{url}            % simple URL typesetting
\urlstyle{same}
\usepackage{booktabs}       % professional-quality tables
\usepackage{amsfonts}       % blackboard math symbols
\usepackage{amsthm}
\usepackage{amsmath}
\usepackage{nicefrac}       % compact symbols for 1/2, etc.
\usepackage{microtype}      % microtypography
\usepackage{lipsum}
\usepackage[sort,numbers]{natbib} %bibliography settings
\bibliographystyle{plainnat}
\usepackage[title]{appendix}
\usepackage{listings}      % For including code.

% For theorems/definitions etc
\newtheorem{theorem}{Theorem}
\newtheorem{lemma}{Lemma}
\newtheorem{proposition}{Proposition}
\newtheorem{corollary}{Corollary}
\theoremstyle{definition}
\newtheorem{definition}{Definition}
\newtheorem{remark}{Remark}

%%%%%%%%%%%%%%%%%%%%%%%%%%%%%%%%
%
% page layout specifications - 
%
%%%%%%%%%%%%%%%%%%%%%%%%%%%%%%%%

% for 25 mm margins all around (left, right and bottom to text, top to header)

\newlength{\extraspace}
\setlength{\extraspace}{\lineskip}

%vertical margins
%==================
% page height is: 297mm
\voffset -1in
\topmargin 25mm
% from header down height is: 272mm
\headheight 4.3mm %just a little higher than 12pt = 12x0.3527 = 4.2324
\headsep 4mm
% from text down height is: 263.7mm
\textheight 238.7mm % compensate for 25mm bottom margin
\addtolength{\headsep}{2.5mm} %removes annoying lineskip below last line of text on page

%horizontal margins
%==================
% page width is: 210mm
\hoffset -1in
\marginparwidth 0mm
\marginparsep 0mm
\evensidemargin 25mm
\oddsidemargin 25mm
\textwidth 160mm


% Set paragraph behaviour
%========================
\parindent 0mm
\parskip 2mm plus 1pt

%%%%%%%%%%%%%%%%%%%%%%%%%%%%%%%%%%%%%%%%%%%%%%%%%%
%%%%%%%%%%%%%%%%%%%%% %%%%%%%%%%%%%%%%%%%%%%%%%
%%%%%%%%%%%%%%%%%%%%% YOU MAY NOW MAKE CHANGES! 
%%%%%%%%%%%%%%%%%%%%% %%%%%%%%%%%%%%%%%%%%%%%%% 
%%%%%%%%%%%%%%%%%%%%%%%%%%%%%%%%%%%%%%%%%%%%%%%%%%

\title{A template for Applied Mathematics Honours project}

\author{A.S.\ tu Dent\\ % Put your name here. Note the \ after . prevents a large white space.
\hspace*{0cm}\\
Supervisor: Dr S.\ Omebody}
\date{} %Do not add date

\begin{document}
\maketitle % Leave this.

\begin{abstract}
Type abstract here.
\lipsum[1]
\end{abstract}

%%---------------------------------------------------------------------------------------------------------------------------------------------%%
%% Introduction section containing problem statement                                                                                                           %%
%%---------------------------------------------------------------------------------------------------------------------------------------------%%
\section{Introduction}
\lipsum[2]
\lipsum[3]

%%---------------------------------------------------------------------------------------------------------------------------------------------%%
%% Examples of heading levels. Add sections as necessary                                                                                                    %%
%%---------------------------------------------------------------------------------------------------------------------------------------------%%
\section{Headings: first level}
\label{sec:headings}

\lipsum[4] See Section \ref{sec:headings}.

\subsection{Headings: second level}
\lipsum[5] 

\subsubsection{Headings: third level}
\lipsum[6]

\paragraph{Paragraph}
\lipsum[7]

\newpage % Try to avoid using newpage. Latex will usually do its best to make things pretty autmomatically.

\section{Equations \& Theorems}

\subsection{Equations}

Short equations cane be written inline, such as $e^{i\pi} = -1$. Longer equations should have their own line. However, if equations are not part of sentences, the probability that it will be frowned upon is
\begin{equation}\label{eqn:frown}
	\xi _{ij}(t)=P(x_{t}=i,x_{t+1}=j|y,v,w;\theta)= {\frac {\alpha _{i}(t)a^{w_t}_{ij}\beta _{j}(t+1)b^{v_{t+1}}_{j}(y_{t+1})}{\sum _{i=1}^{N} \sum _{j=1}^{N} \alpha _{i}(t)a^{w_t}_{ij}\beta _{j}(t+1)b^{v_{t+1}}_{j}(y_{t+1})}}.
\end{equation}
Use \verb+\label+ and \verb+\ref+ to automatically refer to equations and figures, like Equation~(\ref{eqn:frown}), Figure~\ref{fig:fig1}, and Table~\ref{tab:table}. % Pro tip. The ~ between the word and the reference makes sure it keeps a nice space and does not drop to a new line.

Using \verb+\begin{equation}...\end{equation}+ will automatically number an equation. Using \verb+\[...\]+ or \verb+$$...$$+ will not.

\subsection{Definitions and Theorems}
The {\tt amsthm} package provides a theorem environment. Theorems, definitions, etc, will be numbered automatically, like equations.

\begin{definition}[Lipschitz continuity]
	A function $f:\mathbb{R}\to\mathbb{R}$ is \emph{Lipschitz continuous} if there exists a constant $L\ge 0$ such that
	\begin{equation}
	|f(x)-f(y)| \le L|x-y| \quad \text{for all } x,y\in\mathbb{R}.
	\end{equation}
\end{definition}

\begin{theorem}
	Let $f:\mathbb{R}\to\mathbb{R}$ be Lipschitz continuous with constant $L$. Then $f$ is continuous on $\mathbb{R}$.
\end{theorem}

\begin{proof}
	Let $x\in\mathbb{R}$ and let $\varepsilon>0$ be given. Choose $\delta=\varepsilon/L$ if $L>0$ (and any $\delta>0$ if $L=0$).
	If $|x-y|<\delta$, then by the Lipschitz condition,
	\begin{equation}
	|f(x)-f(y)| \le L|x-y| < L\delta = \varepsilon.
	\end{equation}
	Hence $f$ is continuous at $x$. Since $x$ was arbitrary, $f$ is continuous on $\mathbb{R}$.
\end{proof}

%%---------------------------------------------------------------------------------------------------------------------------------------------%%
%% Examples of citations, figures, tables abd references. Add sections as necessary                                                          %%
%%---------------------------------------------------------------------------------------------------------------------------------------------%%
\section{Examples of citations, figures, tables, references}
\label{sec:others}
\lipsum[8] \cite{dabney2009linear,kour2014fast} and see \citet[page 9]{hadash2018estimate}.

The documentation for \verb+natbib+ may be found at
\begin{center}
  \url{http://mirrors.ctan.org/macros/latex/contrib/natbib/natnotes.pdf}
\end{center}
Of note is the command \verb+\citet+, which produces citations
appropriate for use in inline text.  For example,
\begin{verbatim}
   \citet{hasselmo} investigated\dots
\end{verbatim}
produces
\begin{quote}
  Hasselmo, et al.\ (1995) investigated\dots
\end{quote}

\begin{center}
  \url{https://www.ctan.org/pkg/booktabs}
\end{center}


\subsection{Figures}

See Figure \ref{fig:fig1}. Also, here is how you add footnotes.\footnote{The footnote should go after the sentence ends, without a space.} 

\begin{figure}[t!]
  \centering
  %\includegraphics{}
  \fbox{\rule[-.5cm]{4cm}{4cm} \rule[-.5cm]{4cm}{0cm}}
  \caption{Sample figure caption.}
  \label{fig:fig1}
\end{figure}

You have {\em some} control over where a figure appears by using \verb+\begin{figure}[h]+ (here) or  \verb+\begin{figure}[t]+ (top of page), but results may vary.

\subsection{Tables}
\lipsum[12]
See awesome Table~\ref{tab:table}.

\begin{table}
 \caption{Sample table title}
  \centering
  \begin{tabular}{lll}
    \toprule
    \multicolumn{2}{c}{Part}                   \\
    \cmidrule(r){1-2}
    Name     & Description     & Size ($\mu$m) \\
    \midrule
    Dendrite & Input terminal  & $\sim$100     \\
    Axon     & Output terminal & $\sim$10      \\
    Soma     & Cell body       & up to $10^6$  \\
    \bottomrule
  \end{tabular}
  \label{tab:table}
\end{table}

\subsection{Lists}
\begin{itemize}
\item Lorem ipsum dolor sit amet.
\item consectetur adipiscing elit. 
\item Aliquam dignissim blandit est, in dictum tortor gravida eget. In ac rutrum magna.
\end{itemize}
Alternatively, 
\begin{enumerate}
	\item Lorem ipsum dolor sit amet.
	\item consectetur adipiscing elit. 
	\item Aliquam dignissim blandit est, in dictum tortor gravida eget. In ac rutrum magna.
\end{enumerate}

%%---------------------------------------------------------------------------------------------------------------------------------------------%%
%% Future work                                                                                                                                                                        %%
%%---------------------------------------------------------------------------------------------------------------------------------------------%%
\section{Future work}

\lipsum[13]

%%---------------------------------------------------------------------------------------------------------------------------------------------%%
%% Conclusion                                                                                                                                                                          %%
%%---------------------------------------------------------------------------------------------------------------------------------------------%%
\section{Conclusion}

\lipsum[14]

\section{AI usage declaration (example)}

Sections 2 and 4 were initially drafted with assistance from OpenAI's GPT-4. These 
sections were subsequently substantially revised, expanded, and edited to 
incorporate my own analysis and ensure alignment with my research findings. All 
statistical interpretations, conclusions, and key arguments represent my 
independent work and critical assessment of the data.

%%---------------------------------------------------------------------------------------------------------------------------------------------%%
%% References                                                                                                                                                                        %%
%%---------------------------------------------------------------------------------------------------------------------------------------------%%

%% There are different styles available.
%\bibliographystyle{abbrv}  
%\bibliographystyle{plain} 
%\bibliographystyle{alpha}
%\bibliographystyle{siam}
\bibliography{references}



%%---------------------------------------------------------------------------------------------------------------------------------------------%%
%% Appendices                                                                                                                                                                         %%
%%---------------------------------------------------------------------------------------------------------------------------------------------%%
%% Please begin each appendix on a new page

\newpage

\begin{appendices}
\section{Project plan}

\newpage
\section{Something else}
\end{appendices}
\end{document}
